\begin{theorem}[Conditional group-action $Z$-gauge reformulation of the symmetry corollary]%
    \label{cor:cor_4_1_physical_symmetry_zgauge}
    \lean{MPSTensor.cor_4_1_physical_symmetry_zgauge}
    \leanok
    \uses{def:irreducible_form, def:z_gauge_equiv, def:on_site_symmetric,
        def:twisted_tensor, thm:periodic_ft}
    Let $A$ be an MPS tensor in irreducible form~II and let
    $U : G \to \GL_d(\C)$ be a representation of a group $G$ on the
    physical leg, acting unitarily ($U(g)\, U(g)^{\dagger} = \Id$).
    Assume moreover the periodic equal-case Fundamental Theorem of MPS
    (Theorem~\ref{thm:periodic_ft}) as an explicit hypothesis on the
    pair $(d, D)$, expressed by the hypothesis
    \lean{MPSTensor.PeriodicEqualCaseFT}\,$d\,D$.
    If $A$ is on-site symmetric under~$U$ --- that is, the twisted tensor
    $\widetilde{A}_g$ has the same MPV family as~$A$ for every $g \in G$
    --- then for each $g \in G$ there exists a positive period $m_g$ and
    a $\Z_{m_g}$-gauge equivalence between~$A$ and~$\widetilde{A}_g$.

    Concretely there are matrices $Z_g$ and~$Y_g$, with $Z_g^{m_g} = \Id$
    and $[A^i, Z_g] = 0$, satisfying, for every physical index~$i$,
    \[
        Z_g\, A^i = Y_g\, \widetilde{A}_g^{\,i}\, Y_g^{-1}
    \]

    % Source comparison: the source corollary in
    % \cite[lines~834--845]{DeLasCuevas2017Irreducible} is stated for a single
    % local unitary. It concludes that, for a local unitary $u$, there are a
    % diagonal unitary $Z$ and a virtual unitary $Y$ satisfying, for every
    % physical index $i'$,
    % \[
    %     \sum_i u^{i',i} A^i = Z\,Y\,A^{i'}Y^{\dagger}.
    % \]
    % The statement here is a pointwise group-action reformulation using the
    % periodic equal-case Fundamental Theorem as an explicit hypothesis. It does
    % not assert the diagonal normalization of $Z_g$ or the exact displayed
    % source equation.
\end{theorem}

\begin{proof}\leanok
    \uses{def:rotate_physical, thm:periodic_ft,
        thm:irreducible_form_rotate_physical, cor:symmetry_periodic_assembly}
    The twisted tensor $\widetilde{A}_g$ coincides with the physical-index
    rotation $A_{U(g)}$ of Definition~\ref{def:rotate_physical}.
    By Theorem~\ref{thm:irreducible_form_rotate_physical} the rotated
    tensor remains in irreducible form~II, and the on-site symmetry
    hypothesis supplies the equality $\mathcal{V}(A) = \mathcal{V}(A_{U(g)})$.
    Apply the periodic equal-case Fundamental Theorem
    (Theorem~\ref{thm:periodic_ft}) to obtain the asserted
    $\Z_{m_g}$-gauge equivalence.
\end{proof}

\begin{definition}[Conditional periodic projective-rigidity hypothesis]%
    \label{def:periodic_projective_rigidity}
    \lean{MPSTensor.PeriodicProjectiveRigidity}
    \leanok
    \uses{def:twisted_tensor, def:on_site_symmetric}
    Let $A$ be an MPS tensor with physical dimension $d$ and bond dimension $D$,
    and assume it is on-site symmetric under a group action
    $U : G \to \GL_d(\C)$. The tensor $A$ satisfies the
    \emph{periodic projective-rigidity hypothesis} if one may choose a
    family of virtual gauges $Y : G \to \GL_D(\C)$ and a scalar
    $2$-cochain $u : G \times G \to \C^{\times}$ such that

    \begin{enumerate}
        \item every $Y_g$ realizes the pointwise $Z$-gauge intertwining
              modeled on the symmetry corollary with the twisted tensor
              $\widetilde{A}_g$ for some accompanying
              $Z_g$ (with $Z_g^{m_g} = \Id$ and $[A^i, Z_g] = 0$); and
        \item the family obeys the projective multiplication law on the
              bond space: $Y_h\, Y_g = u(g, h)\, Y_{gh}$.
    \end{enumerate}

    The source supplies the corresponding intertwining for a single local unitary.
    The group-indexed coherent choice and the projective multiplication law are
    additional assumptions: after proving the symmetry corollary, the source
    paper says that its consequences will be explored elsewhere. The injective
    case reduces this coherence to scalar uniqueness of gauges
    (cf.~\S\ref{sec:injective_symmetry_cocycle}); in the genuinely periodic
    setting there is a non-trivial $\Z_m$-gauge ambiguity, so coherence of the
    factor system $u$ is an independent analytic assumption.
\end{definition}

\begin{theorem}[Conditional projective representation from symmetry gauges]%
    \label{thm:cor_4_1_projective_rep}
    \lean{MPSTensor.cor_4_1_projective_rep}
    \leanok
    \uses{def:periodic_projective_rigidity}
    Assume the periodic projective-rigidity hypothesis
    (Definition~\ref{def:periodic_projective_rigidity}) for the tensor $A$ and
    the physical action $U$. Then the virtual gauge family supplied by that
    hypothesis defines a projective representation of $G$ on the bond space:
    there exist a scalar $2$-cochain
    $\omega : G \times G \to \C^{\times}$ and virtual gauges
    $X : G \to \GL_D(\C)$ satisfying
    \[
        X(g)\, X(h) =  \omega(g, h)\, X(g h),
    \]
    while for each $g$ the matrix $X(g^{-1})$ still realizes the pointwise
    $\Z_{m_g}$-intertwining, encoded in the rigidity hypothesis, between $A$
    and the twisted tensor $\widetilde{A}_g$.
\end{theorem}

\begin{proof}\leanok
    \uses{def:periodic_projective_rigidity}
    Unpack the rigidity hypothesis to extract the family $(Y, u)$ together
    with the pointwise $Z$-gauge intertwiners. Set
    $X(g) := Y(g^{-1})$ and $\omega(g, h) := u(h^{-1}, g^{-1})$. The
    projective multiplication law $X(g)\, X(h) = \omega(g, h)\, X(gh)$
    is then exactly the $(h^{-1}, g^{-1})$-instance of the factor-system
    identity, using $(gh)^{-1} = h^{-1} g^{-1}$. The intertwining
    relation between $A$ and $\widetilde{A}_g$ transports unchanged.
\end{proof}

\begin{theorem}[Cocycle identity for the upgraded factor system]%
    \label{thm:cor_4_1_projective_rep_cocycle}
    \lean{MPSTensor.cor_4_1_projective_rep_cocycle}
    \leanok
    \uses{thm:cor_4_1_projective_rep}
    If the bond dimension $D$ is positive, the factor system $\omega$
    produced by Theorem~\ref{thm:cor_4_1_projective_rep} satisfies the
    multiplicative $2$-cocycle identity
    $\omega(g, h)\, \omega(g h, k) = \omega(g, h k)\, \omega(h, k)$.
\end{theorem}

\begin{proof}\leanok
    \uses{thm:cor_4_1_projective_rep}
    Associativity of matrix multiplication applied to
    $X(g)\, X(h)\, X(k)$, combined with the projective multiplication
    law, gives
    \begin{equation}\label{eq:pft_assoc_scalar_multiple}
        \omega(g, h)\, \omega(g h, k)\, X((g h) k)
        =
        \omega(g, h k)\, \omega(h, k)\, X(g (h k)).
    \end{equation}
    Since $(g h) k = g (h k)$, both sides of
    (\ref{eq:pft_assoc_scalar_multiple}) are scalar multiples of the
    same matrix $X((g h) k)$. The matrix lies in $\GL_D(\C)$ and is
    therefore invertible, but the scalar cancellation leading from an
    equality of scalar multiples of a matrix to equality of the scalar
    coefficients additionally requires $D > 0$ so that the underlying
    index set is inhabited. Cancelling then yields
    $\omega(g, h)\, \omega(g h, k) = \omega(g, h k)\, \omega(h, k)$
    in $\C^{\times}$, hence in $\mathrm{Units}(\C)$.
\end{proof}

\begin{remark}[External cohomological interpretation]%
    \label{rem:cor_4_1_projective_rep_spt}
    \uses{thm:cor_4_1_projective_rep_cocycle}
    In the one-dimensional SPT literature, a $U(1)$-valued factor system
    of a symmetric MPS gives a cohomology class in $H^{2}(G,U(1))$
    (cf.~Chen--Gu--Wen and P\'erez-Garc\'ia et al.). This interpretation
    is external to \cite{DeLasCuevas2017Irreducible}: Section~4.2 proves
    the $Z$-gauge symmetry corollary and leaves its consequences to
    later work. The result established here is only the 2-cocycle
    identity of Theorem~\ref{thm:cor_4_1_projective_rep_cocycle}, under
    the conditional rigidity hypothesis of
    Definition~\ref{def:periodic_projective_rigidity}.
\end{remark}

\begin{theorem}[Conditional projective representation and cocycle from symmetry]%
    \label{thm:projectiveRep_of_symmetry}
    \lean{MPSTensor.projectiveRep_of_symmetry}
    \leanok
    \uses{thm:cor_4_1_projective_rep, thm:cor_4_1_projective_rep_cocycle}
    Theorems~\ref{thm:cor_4_1_projective_rep}
    and~\ref{thm:cor_4_1_projective_rep_cocycle} combine into a
    single statement: from the periodic projective-rigidity
    hypothesis (Definition~\ref{def:periodic_projective_rigidity}), one
    obtains a projective representation of $G$ on the bond space whose
    factor system is a genuine $2$-cocycle whenever $D > 0$.
\end{theorem}
